\bibitem{R_A17}
migodam，A17最终科学closure（用户项目原始证据），commit \texttt{df61a1b5}，2026。
\href{https://github.com/migodam/a17-maxwell-certified-exchange/blob/df61a1b533682a883cdaf65259b4528c354bd832/closure_final/reports/A17_CLOSURE_REPORT_ZH.md}{固定版本原文}。
入口：\href{https://github.com/migodam/a17-maxwell-certified-exchange/blob/df61a1b533682a883cdaf65259b4528c354bd832/closure_final/START_HERE.md}{固定版本原文}。

\bibitem{R_A17core}
migodam，A17正式理论：CORE\_THEORY与CONDITIONAL\_EXCHANGE\_THEORY，同上commit。
\href{https://github.com/migodam/a17-maxwell-certified-exchange/blob/df61a1b533682a883cdaf65259b4528c354bd832/closure_final/manuscript/CORE_THEORY.tex}{固定版本原文}；
\href{https://github.com/migodam/a17-maxwell-certified-exchange/blob/df61a1b533682a883cdaf65259b4528c354bd832/closure_final/manuscript/CONDITIONAL_EXCHANGE_THEORY.tex}{固定版本原文}。

\bibitem{R_A17eff}
migodam，A17定向核验效率收口，同上commit。
\href{https://github.com/migodam/a17-maxwell-certified-exchange/blob/df61a1b533682a883cdaf65259b4528c354bd832/efficiency_followup/A17_EFFICIENT_VERIFICATION_FINAL.md}{固定版本原文}。

\bibitem{R_A26}
migodam，A26最终科学裁决，commit \texttt{2603f3bd}，2026。
\href{https://github.com/migodam/a26-om-feedback-prior-completion/blob/2603f3bd16c48e48587fedc4363975779def17f5/A26_FINAL_DECISION.md}{固定版本原文}。

\bibitem{R_A30}
migodam，A30 E0/E1材料覆盖、测量选择与合法补偿，commit \texttt{717d3d26}，2026。
\href{https://github.com/migodam/a30-maxwell-observable-material-manifold/blob/717d3d26a290248bd2b9aeb55f3d03694b407b84/A30_E0_E1_DECISION.md}{固定版本原文}。

\bibitem{R_B1}
migodam，B1总入口及B1--C1 Joint P2科学裁决，commit \texttt{064d0aae}，2026。
\href{https://github.com/migodam/b1-probabilistic-maxwell-imaging/blob/064d0aaed9772cadd573866eec84f159895b317b/START_HERE.md}{固定版本原文}；
\href{https://github.com/migodam/b1-probabilistic-maxwell-imaging/blob/064d0aaed9772cadd573866eec84f159895b317b/research/joint_p2/GATE_DECISION.md}{固定版本原文}。

\bibitem{R_C1}
migodam，C1 P2子模块科研裁决，commit \texttt{83056881}，2026。
\href{https://github.com/migodam/c1-physics-state-generative-imaging/blob/83056881e89329212b16f5be47a4d0dbc4465943/p2/C1_P2_DECISION.md}{固定版本原文}。

\bibitem{R_P2}
migodam，Research P2：存在有限噪声材料自由度，但尚无独立的软谱生成算法优势，B1 commit \texttt{064d0aae}，2026-10-19。
\href{https://github.com/migodam/b1-probabilistic-maxwell-imaging/blob/064d0aaed9772cadd573866eec84f159895b317b/research/p2_material_fiber_20261019/RESEARCH_REPORT_ZH.md}{固定版本原文}。

\bibitem{R_F3}
migodam，Research F3：商坐标解决了驻点障碍，但没有支持高阶Explorer，同上B1 commit，2026-10-17。
\href{https://github.com/migodam/b1-probabilistic-maxwell-imaging/blob/064d0aaed9772cadd573866eec84f159895b317b/research/research_f3_quotient_exploration/RESEARCH_REPORT_ZH.md}{固定版本原文}。

\bibitem{R_A17X}
migodam，Research A17-X原Pro报告与A17-X/F3竞争审计，同上B1 commit，2026。
\href{https://github.com/migodam/b1-probabilistic-maxwell-imaging/blob/064d0aaed9772cadd573866eec84f159895b317b/research/a17x_f3_competition_audit/pro_source/report/RESEARCH_A17X_FULL_REPORT_ZH.md}{固定版本原文}；
\href{https://github.com/migodam/b1-probabilistic-maxwell-imaging/blob/064d0aaed9772cadd573866eec84f159895b317b/research/a17x_f3_competition_audit/RESEARCH_REPORT_ZH.md}{固定版本原文}。
原Pro为理论研究资料，实施审计为独立证据类别。

\bibitem{R_F}
migodam，Research F原Pro理论报告与材料几何实施报告，同上B1 commit，2026。
\href{https://github.com/migodam/b1-probabilistic-maxwell-imaging/blob/064d0aaed9772cadd573866eec84f159895b317b/research/research_f_intrinsic_transport/pro_source/report/RESEARCH_REPORT_ZH.md}{固定版本原文}；
\href{https://github.com/migodam/b1-probabilistic-maxwell-imaging/blob/064d0aaed9772cadd573866eec84f159895b317b/research/research_f_intrinsic_transport/RESEARCH_REPORT_ZH.md}{固定版本原文}。

\bibitem{R_B}
migodam，Research B：补偿感知主动Maxwell测量的实际实验，同上B1 commit，2026-10-15。
\href{https://github.com/migodam/b1-probabilistic-maxwell-imaging/blob/064d0aaed9772cadd573866eec84f159895b317b/research/research_b/RESEARCH_REPORT_ZH.md}{固定版本原文}。

\bibitem{R_A}
migodam，Research A：从单张重建转向顺序发现相容材料世界，同上B1 commit，2026-10-15。
\href{https://github.com/migodam/b1-probabilistic-maxwell-imaging/blob/064d0aaed9772cadd573866eec84f159895b317b/research/research_a_20261015/RESEARCH_REPORT_ZH.md}{固定版本原文}。

\bibitem{R_T}
migodam，Research T：能补偿出新孔洞，但未改善成像，同上B1 commit，2026-10-15。
\href{https://github.com/migodam/b1-probabilistic-maxwell-imaging/blob/064d0aaed9772cadd573866eec84f159895b317b/research/research_t_20261015/RESEARCH_REPORT_ZH.md}{固定版本原文}。

\bibitem{R_G}
migodam，Research G：统一材料空间中的下一解搜索，同上B1 commit，2026-10-15。
\href{https://github.com/migodam/b1-probabilistic-maxwell-imaging/blob/064d0aaed9772cadd573866eec84f159895b317b/research/research_g_20261015/RESEARCH_G_REPORT_ZH.md}{固定版本原文}。

\bibitem{R_LXZ}
migodam，LXZ局部--全局Maxwell多解实施研究，同上B1 commit，2026-10-16。
\href{https://github.com/migodam/b1-probabilistic-maxwell-imaging/blob/064d0aaed9772cadd573866eec84f159895b317b/research/research_lxz_20261016/RESEARCH_REPORT_ZH.md}{固定版本原文}。

